\noindent In this section, we discuss algebraic properties of $\mu_n, n \in \{13,17,21\}$, which are essential for showing that a tiling arising from $\sigma_n$ has ILC. 

\begin{lem}
\label{lem:mu_n_properties}
For each $n \in \{13,17,21\}$, the inflation factor $\mu_n$ is a non-Pisot algebraic integer, and the minimal polynomial $g_n(x)$ of $\mu_{n}$ has no repeated roots.
\end{lem}

\begin{proof}
The representations for $\mu_n$ and $\mu_n s_{n,k}$ as sums of prototile edge lengths, given in Equations~\ref{eq:legrep} and~\ref{eq:baserep}, naturally give rise to the matrix equation  $\mathbf{z}^\top M_n = \mu_n \mathbf{z}^\top,$ where 
\[
\mathbf{z} = 
\begin{bmatrix}
s_{n,0}=1 \\ s_{n,1} \\ s_{n,2} \\ \vdots \\
s_{n,(n-3)/2} \\ s_{n,(n-1)/2}
\end{bmatrix} \  \mbox{and} \ 
M_n = 
 \begin{bmatrix} 1 & 1 & 1 & \cdots &1  & 1 & 1 \\ 
0 & 1 & 1 & \cdots & 1 & 1 & 0 \\ 
1 & 1 & 1 & \cdots & 1 & 0 & 0 \\ 
0 & 1 & 1 & \cdots & 0 & 0 & 0 \\ 
\vdots & \vdots & \vdots & \reflectbox{$\ddots$} & \vdots & \vdots & \vdots \\ 
0 & 1 & 0 & \cdots & 0 & 0 & 0 
\\ 1 & 0 & 0 & \cdots & 0 & 0 & 0 \\ \end{bmatrix}.
\]

Each component of the row vector $\mathbf{z}^\top M_n$ corresponds to a representation of $\mu_n s_{n,k}$, $k=0, 1, \ldots, (n-1)/2$, in terms of the edge lengths $s_{n,j}$. For example, the first entry encodes the representation  $s_{n,0} + s_{n,2} + \cdots + s_{n,(n-1)/2}$ of $\mu_n$ as stated in Equation~\ref{eq:legrep}, while the second entry corresponds to  $s_{n,0} + s_{n,1} + \cdots + s_{n,(n-3)/2} = \mu_n s_{n,1}$ following Equation~\ref{eq:baserep}.  

This confirms that $\mu_n$ is an eigenvalue of the integer matrix $M_n$, with left eigenvector $\mathbf{z}^\top$. Hence, $\mu_{n}$ is an algebraic integer, and its minimal polynomial $g_{n}(x)$ divides the characteristic polynomial $p_n(x)$ of $M_n$.

To determine $g_n(x)$ explicitly, we compute the characteristic polynomial  $p_n(x) = \det(M_n - xI)$, where $I$ denotes the identity matrix of size $(n+1)/2$. For each $n =13,17,21$, we have

\begin{equation}
p_n(x)=-\left(x-1 \right) \left(x+\dfrac{1}{s_{n,1}}\right)\left(x-\dfrac{1}{s_{n,2}} \right) \cdots \left(x+\dfrac{1}{s_{n,(n-3)/2}} \right)\left(x-\dfrac{1}{s_{n,(n-1)/2}} \right),
\end{equation}

\noindent with $\mu_{n} = \dfrac{1}{s_{n,(n-1)/2}}$. Thus, the eigenvalues are the positive or negative reciprocals of the prototile edge lengths. Note that the $(n+1)/2$ roots are all distinct, so each eigenvalue is simple. So, $g_{n}(x)$ has no repeated roots. For $n = 13$ and $n = 17$, the minimal polynomial is $g_{n}(x) = \frac{p_{n}(x)}{-(x - 1)},$ while for $n = 21$ we have  $g_{21}(x) = \frac{p_{21}(x)}{-\,(x^2 - 1)\,(x^3 + 2x^2 - x - 1)}.$ These polynomials are listed in Table~\ref{tab:minpoly}, together with the algebraic conjugate $\mu_{n,2}$ of $\mu_n$ whose modulus exceeds $1$, confirming that $\mu_n$ is non-Pisot.
\newpage
\begin{center}
\captionof{table}{The minimal polynomial of $\mu_n$ and its conjugate $\mu_{n,2}$ with $|\mu_{n,2}|>1.$}
\renewcommand{\arraystretch}{1.5}
\label{tab:minpoly}
\begin{tabular}{| c | c | c |}
\hline
$n$ & Minimal Polynomial $g_n (x)$ &	$\mu_{n,2}$ \\
\hline
$13$ & $x^6-3x^5-6x^4+4x^3+5x^2-x-1$ & $-\dfrac{1}{s_{13,5}} \approx -1.41002$ \\ \hline
$17$ & $x^8-4x^7-10x^6+10x^5+15x^4-6x^3-7x^2+x+1$ & $-\dfrac{1}{s_{17,7}} \approx -1.82706$ \\ \hline
$21$ & $x^6-8x^5+8x^4+6x^3-6x^2-x+1$ & $\dfrac{1}{s_{21,8}} \approx 1.36858$ \\
\hline
\end{tabular}
\end{center}
\end{proof}

The non-Pisot property of the inflation factor will become crucial later in the proof of ILC, where it will allow us to construct an infinite sequence of distinct numbers. These numbers will come from lengths of certain segments in the tiling, and our first step toward that goal is to consider the algebraic setting in which these lengths live.

Since $\mu_n$ is an algebraic integer, it generates the ring $\mathbb{Z}[\mu_n]$, which is a free $\mathbb{Z}$-module of rank $d_n = \deg g_n(x)$. 
This module can be naturally identified with $\mathbb{Z}^{d_n}$ via the standard (monomial) basis $B = \{1, \mu_n, \mu_n^2, \dots, \mu_n^{d_n - 1}\}.$


For each $s \in \mathbb{Z}[\mu_n]$, let  
$$
s = a_0 + a_1 \mu_n + a_2 \mu_n^2 + \cdots + a_{d_n-1} \mu_n^{d_n-1}, 
\quad a_i \in \mathbb{Z},
$$  
be its unique representation in terms of the basis $B$. The mapping  
$$
s \longmapsto (a_0, a_1, \dots, a_{d_n-1})^\top
$$  
is the \textit{coordinate isomorphism} with respect to $B$, and $(a_0, a_1, \dots, a_{d_n-1})^\top$ is called the \textit{coordinate vector} of $s$ with respect to $B$. For simplicity, whenever we refer to the coordinate vector of an element of $\mathbb{Z}[\mu_n]$, we mean with respect to the standard basis $B$.

In the next lemma, we establish that $s_{n,j} \in \mathbb{Z}[\mu_n]$ and then compute their coordinate vectors.

\begin{lem}
For each \( n \in \{13, 17, 21\} \) and \( j \in I_n=\{0,1,2,\ldots,(n-1)/2\} \), we have \( s_{n,j} \in \mathbb{Z}[\mu_n] \).
\label{lem:s_n,j}
\end{lem}

\begin{proof}
Let \(\mathbf{o}\) denote the zero vector of dimension \((n+1)/2\). Since \(\mathbf{z}\) is an eigenvector of \(M_n\) with eigenvalue \(\mu_n\) by Lemma~\ref{lem:mu_n_properties}, it satisfies the homogeneous system
\[
(M_n^\top - \mu_n I)\,\mathbf{x} = \mathbf{o}.
\]

Solving this system by Gaussian elimination expresses each \(s_{n,j}\) as a \(\mathbb{Q}\)-linear combination of powers of \(\mu_n\). Reducing these powers using the relation \(g_n(\mu_n) = 0\) shows that each \(s_{n,j}\) can be written in the standard basis \(B\) of \(\mathbb{Z}[\mu_n]\). The corresponding coordinate vectors are listed in Table~\ref{tab:vectorrep}.
\end{proof}

As a final preparatory step,  we recall how multiplication in $\mathbb{Z}[\mu_n]$ by $\mu_n$ corresponds to $\mathbb{Z}^{d_n}$. Let $G_n$ denote the companion matrix of $g_n(x)$. For $s \in \mathbb{Z}[\mu_n]$ with coordinate vector $(a_0, a_1, \dots, a_{d_n - 1})^\top$, the vector $G_n(a_0, a_1, \dots, a_{d_n - 1})^\top$ is precisely the coordinate vector of $\mu_n s$. 

By Lemma~\ref{lem:mu_n_properties}, the polynomial $g_n(x)$ has no repeated roots. Consequently, the companion matrix $G_n$ is diagonalizable over the splitting field of $g_n$. Thus there exists an invertible matrix $A$ such that
\begin{equation}
\label{eq:Gtilde}
A^{-1} G_n A = \widetilde{G}_n =
\begin{bmatrix}
\mu_n & 0 & 0 & \cdots & 0 \\
0 & \mu_{n,2} & 0 & \cdots & 0 \\
0 & 0 & \mu_{n,3} & \cdots & 0 \\
\vdots & \vdots & \vdots & \ddots & \vdots \\
0 & 0 & 0 & \cdots & \mu_{n,d_n}
\end{bmatrix},
\end{equation}
where the $i$th column of $A$ is a right eigenvector and the $i$th row of $A^{-1}$ is a left eigenvector of $G_n$ corresponding to the eigenvalue $\mu_{n,i}$, for $i = 1, 2, \ldots, d_n$, with $\mu_{n,1} = \mu_n$. We also note that $\mu_{n,2}$ is the algebraic conjugate of $\mu_n$ with modulus greater than $1$, as given in Table~\ref{tab:minpoly}.

In this work, we choose $A^{-1}$ so that its second row is
\[
v_{n,2} =
\left( \frac{1}{\mu_{n,2}^{d_n - 1}},\ \frac{1}{\mu_{n,2}^{d_n - 2}},\ \ldots,\ \frac{1}{\mu_{n,2}},\ 1 \right).
\]
As will be needed later, we also include \( v_{n,2} y_{n,j} \) in Table~\ref{tab:vectorrep}.
\newpage

\begin{table}[htbp]
\centering
\caption{The coordinate vector $y_{n,j}$ of $s_{n,j}$ with respect to the standard basis $B$; and $v_{n,2} {y}_{n,j}$. }
\label{tab:vectorrep}
\begin{tabular}{|l|l|l|}
\hline
\( s_{n,j} \) & Coordinate vector $y_{n,j}$ & \( v_{n,2} {y}_{n,j} \) \\
\hline
\( s_{13,0} \) & \( (1,0,0,0,0,0)^\top \) & -0.179421119170724 \\ \hline
\( s_{13,1} \) & \( (-4,1,10,-3,-4,1)^\top \)& -0.268597272296381 \\ \hline
\( s_{13,2} \) & \( (5,-4,-20,11,11,-3)^\top \) & -0.0432536521449611 \\ \hline
\( s_{13,3} \) & \( (0,5,5,-11,-6,2)^\top \) & 0.203845625239822 \\ \hline
\( s_{13,4} \) & \( (-5,0,16,-5,-8,2)^\top \) & 0.348414935064387 \\ \hline
\( s_{13,5} \) & \( (4,-5,-15,14,10,-3)^\top \) & 0.317739022198404 \\ \hline
\( s_{13,6} \) & \( (-1,5,4,-6,-3,1)^\top \) & 0.127247211393141 \\
\hline
\( s_{17,0} \) & \( (1,0,0,0,0,0,0,0)^\top \) & -0.0147137185440312 \\ \hline
\( s_{17,1} \) & \( (-6,1,21,-5,-20,6,5,-1)^\top \) & -0.0250197112728618 \\ \hline
\( s_{17,2} \) & \( (14,-6,-70,29,85,-33,-24,5)^\top \) & -0.0131169374224207 \\ \hline
\( s_{17,3} \) & \( (-14,14,84,-64,-126,67,41,-9)^\top \) & 0.00271522134344693 \\ \hline
\( s_{17,4} \) & \( (0,-14,-14,56,42,-48,-20,5)^\top \) & 0.0177339928490863 \\ \hline
\( s_{17,5} \) & \( (14,0,-70,14,90,-28,-25,5)^\top \) & 0.02744026786708 \\ \hline
\( s_{17,6} \) & \( (-14,14,78,-70,-117,71,40,-9)^\top \) & 0.0289263790503149 \\ \hline
\( s_{17,7} \) & \( (6,-14,-35,55,56,-45,-21,5)^\top \) & 0.0217471384189736 \\ \hline
\( s_{17,8} \) & \( (-1,7,6,-15,-10,10,4,-1)^\top \) & 0.0080532004236864 \\ 
\hline
\( s_{21,0} \) & \( (1,0,0,0,0,0)^\top \) & 0.208277419875994 \\ \hline
\( s_{21,1} \) & \( (-5,0,14,0,-7,1)^\top \) & 0.305356304999981 \\ \hline
\( s_{21,2} \) & \( (9,0,-35,-1,21,-3)^\top \) & 0.0311291821584927 \\ \hline
\( s_{21,3} \) & \( (-5,1,21,1,-14,2)^\top \) & -0.259717694500226 \\ \hline
\( s_{21,4} \) & \( (-6,-5,27,9,-22,3)^\top \) & -0.411902266359988 \\ \hline
\( s_{21,5} \) & \( (13,10,-59,-26,52,-7)^\top \) & -0.344173760232181 \\ \hline
\( s_{21,6} \) & \( (-10,-9,45,26,-45,6)^\top \) & -0.0926921719855542 \\ \hline
\( s_{21,7} \) & \( (1,0,0,0,0,0)^\top \) & 0.208277419875994 \\ \hline
\( s_{21,8} \) & \( (5,9,-31,-26,38,-5)^\top \) & 0.398048476985537 \\ \hline
\( s_{21,9} \) & \( (-4,-10,24,25,-31,4)^\top \) & 0.375302942390668 \\ \hline
\( s_{21,10} \) & \( (1,6,-6,-8,8,-1)^\top \) & 0.152184571859762 \\
\hline
\end{tabular}
\end{table}



